\clearpage
\section{Reasoning environments}
% example:env-theorem
\begin{theorem}[A useful observation]
\label{theorem:demo}
For $x>0$, the square $x^2$ is positive.
\end{theorem}
% endexample
% example:env-lemma
\begin{lemma}[A useful observation]
\label{lemma:demo}
For $x>0$, the square $x^2$ is positive.
\end{lemma}
% endexample
% example:env-corollary
\begin{corollary}[A useful observation]
\label{corollary:demo}
For $x>0$, the square $x^2$ is positive.
\end{corollary}
% endexample
% example:env-proposition
\begin{proposition}[A useful observation]
\label{proposition:demo}
For $x>0$, the square $x^2$ is positive.
\end{proposition}
% endexample
% example:env-conjecture
\begin{conjecture}[A useful observation]
\label{conjecture:demo}
For $x>0$, the square $x^2$ is positive.
\end{conjecture}
% endexample
% example:env-algorithm
\begin{algorithm}[A useful observation]
\label{algorithm:demo}
For $x>0$, the square $x^2$ is positive.
\end{algorithm}
% endexample
% example:env-definition
\begin{definition}[A useful observation]
\label{definition:demo}
For $x>0$, the square $x^2$ is positive.
\end{definition}
% endexample
% example:env-example
\begin{example}[A useful observation]
\label{example:demo}
For $x>0$, the square $x^2$ is positive.
\end{example}
% endexample
% example:env-remark
\begin{remark}[A useful observation]
\label{remark:demo}
For $x>0$, the square $x^2$ is positive.
\end{remark}
% endexample
\subsection{Review}
% example:env-summary
\begin{summary}[A useful observation]
\label{summary:demo}
For $x>0$, the square $x^2$ is positive.
\end{summary}
% endexample
% example:env-note
\begin{note}[A useful observation]
For $x>0$, the square $x^2$ is positive.
\end{note}
% endexample
% example:env-caveat
\begin{caveat}[A useful observation]
For $x>0$, the square $x^2$ is positive.
\end{caveat}
% endexample
% example:env-warning
\begin{warning}[A useful observation]
For $x>0$, the square $x^2$ is positive.
\end{warning}
% endexample
% example:env-question
\begin{question}[A useful observation]
For $x>0$, the square $x^2$ is positive.
\end{question}
% endexample
% example:env-speculation
\begin{speculation}[A useful observation]
For $x>0$, the square $x^2$ is positive.
\end{speculation}
% endexample
% example:env-proof
\begin{proof}[A useful observation]
For $x>0$, the square $x^2$ is positive.
\end{proof}
% endexample
\clearpage
\section{Displays and lists}
% example:env-formula
\begin{formula}[Area]\label{formula:area}
  A=\pi r^2
\end{formula}
See \cref{formula:area}; the reference is the formula number.
% endexample
% example:env-subparts
\begin{subparts}
  \item Identify the radius.
  \item Calculate the area.
\end{subparts}
% endexample
% example:env-mathtable
\courseenvironmentsetup{table-placement=H,table-top-skip=-1em}
\begin{mathtable}[caption={Squares},label={tab:squares}]{cc}
  \tableheading{Input}{$x$} & \tableheading{Output}{$x^2$}\\
  1 & 1\\ 2 & 4\\ \bottomrule
\end{mathtable}
See \cref{tab:squares}.
% endexample
